\usepackage[utf8]{inputenc}
\usepackage[T1]{fontenc}
\usepackage[provide=*,spanish]{babel}
\usepackage[letterpaper,margin=1.8cm,top=1.9cm,bottom=1.8cm]{geometry}
\usepackage{lmodern,microtype,amsmath,amssymb,booktabs,graphicx,fancyhdr,array,listings,tabularx,tikz,titlesec}
\usepackage[table]{xcolor}
\usepackage[hidelinks]{hyperref}
\definecolor{Tinta}{HTML}{152C45}
\definecolor{Azul}{HTML}{164B70}
\definecolor{Oro}{HTML}{B9955E}
\definecolor{Fondo}{HTML}{F1F5F8}
\definecolor{Gris}{HTML}{556475}
\definecolor{Linea}{HTML}{D3DDE6}
\color{Tinta}
\setlength{\parindent}{0pt}
\setlength{\parskip}{7pt}
\setlength{\emergencystretch}{2em}
\setlength{\headheight}{21pt}
\setlength{\headsep}{15pt}
\pagestyle{fancy}\fancyhf{}
\fancyhead[L]{\sffamily\scriptsize\color{Gris}FUNDAMENTOS DE BASES DE DATOS}
\fancyhead[R]{\sffamily\scriptsize\color{Gris}TAREA 04 / ÁLGEBRA RELACIONAL}
\fancyfoot[L]{\sffamily\scriptsize\color{Gris}UNAM · Facultad de Ciencias}
\fancyfoot[C]{\sffamily\small\thepage}
\fancyfoot[R]{\sffamily\scriptsize\color{Gris}Equipo ÑIPGG}
\renewcommand{\headrule}{\hbox to\headwidth{\color{Linea}\leaders\hrule height .4pt\hfill}}
\renewcommand{\footrulewidth}{0pt}
\titleformat{\section}{\sffamily\Large\bfseries\color{Azul}}{}{0pt}{}[\vspace{2pt}{\color{Oro}\titlerule[.8pt]}]
\titleformat{\subsection}{\sffamily\large\bfseries\color{Azul}}{}{0pt}{}
\titlespacing*{\section}{0pt}{5pt}{12pt}
\titlespacing*{\subsection}{0pt}{13pt}{7pt}
\newcommand{\seccion}[1]{\phantomsection\section*{#1}\addcontentsline{toc}{section}{#1}}
\lstset{basicstyle=\ttfamily\scriptsize\color{Tinta},breaklines=true,columns=fullflexible,frame=single,rulecolor=\color{Linea},backgroundcolor=\color{Fondo},framesep=6pt,xleftmargin=6pt,xrightmargin=6pt,aboveskip=9pt,belowskip=10pt,showstringspaces=false}
\newcommand{\caprel}[3]{\begin{center}\includegraphics[width=\linewidth,height=#2cm,keepaspectratio]{\RutaDiego/evidencias/#1}\par\vspace{4pt}{\sffamily\footnotesize\color{Gris}#3}\end{center}}
\hypersetup{pdftitle={Tarea 04 - Álgebra relacional},pdfauthor={Equipo ÑIPGG},pdfsubject={Fundamentos de Bases de Datos, semestre 2027-I}}

\usepackage{longtable}
\usetikzlibrary{arrows.meta,positioning}
\tikzset{operador/.style={draw=Linea,fill=Fondo,rounded corners=3pt,align=center,font=\sffamily\small,text=Tinta,inner sep=7pt},flujo/.style={-{Latex[length=2mm]},draw=Azul,thick}}
\newcommand{\nulo}{\texttt{null}}

\newcommand{\leftouterjoin}{\mathbin{\ooalign{$\bowtie$\cr\hidewidth\raise.15ex\hbox{\scalebox{1.15}[1]{$\lbrack$}}\kern.7em\hidewidth\cr}}}
\newcommand{\rightouterjoin}{\mathbin{\ooalign{$\bowtie$\cr\hidewidth\kern.7em\raise.15ex\hbox{\scalebox{1.15}[1]{$\rbrack$}}\hidewidth\cr}}}
\lstset{literate={⟕}{{$\leftouterjoin$}}1 {⟖}{{$\rightouterjoin$}}1}
